\bookchapter{Bab 26: Teknologi Ponsel Terkini}

\chapterintro{Setelah membahas “urusan pascamati” ponsel lama, mari kita lihat masa kini. 2024-2026 adalah periode ledakan teknologi ponsel: AI dari cloud masuk ke saku, komunikasi satelit turun dari flagship ke kelas menengah, layar lipat dari “sekadar mencoba” menjadi “pakai sehari-hari”, dan teknologi baterai mengalami terobosan material level yang belum terjadi dalam bertahun-tahun. Bab ini membahas teknologi terdepan yang sedang terjadi, membantu Anda menilai “saatnya ganti ponsel atau belum” dan “ponsel apa yang dipilih”.}

\booksection{Ponsel AI: Hype Pemasaran atau Perubahan Nyata?}

Sejak 2024, hampir semua pabrikan Android mengibarkan bendera “ponsel AI” di produk flagship, dan Apple meluncurkan Apple Intelligence di iOS 18. Tapi ponsel seperti apa yang benar-benar layak disebut “ponsel AI”? Industri telah membentuk tiga kriteria penilaian:

\begin{itemize}
  \item \textbf{On-device large model}: ponsel bisa menjalankan model bahasa berukuran sekitar 7B parameter (seperti versi ringan Qwen, Llama, DeepSeek) secara lokal, tanpa bergantung pada cloud untuk menyelesaikan ringkasan teks, terjemahan, pembuatan copywriting, dll. Cara mengecek: jika fitur AI masih bisa digunakan dalam mode pesawat, berarti modelnya on-device.
  \item \textbf{Daya komputasi NPU (TOPS)}: daya komputasi NPU (Neural Processing Unit) diukur dalam TOPS (triliun operasi per detik). NPU chipset flagship 2024 (seperti Snapdragon 8 Gen3, Dimensity 9300) berada di rentang 40-70 TOPS, generasi baru 2025-2026 umumnya menembus 100 TOPS, ini adalah ambang batas hardware untuk menjalankan large model on-device.
  \item \textbf{Skenario AI Agent}: tidak hanya bisa "Q\& A", tapi juga mengeksekusi rantai tugas lintas App — misalnya "bantu kumpulkan foto rapat kemarin menjadi album dan bagikan ke Zhang San", sistem bisa otomatis memanggil album, file, WeChat dan beberapa App lain untuk menyelesaikannya. Ini adalah tanda inti yang membedakan ponsel AI dari "ponsel dengan asisten AI".
\end{itemize}

\booksection{Perbandingan Strategi AI Pabrikan: Siapa yang Serius, Siapa yang Sekadar Label?}

Strategi AI setiap pabrikan berbeda jelas, saat membeli bisa difokuskan pada:

\begin{itemize}
  \item \textbf{Apple Apple Intelligence (iOS 18+)}: Menekankan integrasi aplikasi level sistem — ringkasan email, prioritas notifikasi, polishing tulisan, AI photo editing album. Terutama on-device model, task kompleus ke cloud ke Private Cloud Compute. Kelebihan adalah perlindungan privasi dan kelancaran sistem, kekurangan adalah skenario bahasa Tiongkok online lebih lambat, sebagian fitur baru bertahap didukung pada 2025.
  \item \textbf{Google Gemini Nano}: tertanam di seri Pixel 8/9, mendukung "Circle to Search", Magic Compose, ringkasan rekaman, dll. Gemini Nano adalah model kecil on-device dalam arti sebenarnya, ekosistem open source ramah, menjadi implementasi referensi ponsel AI Android.
  \item \textbf{: Terintegrasi ke sistem HarmonyOS, asisten Xiaoyi mendapat kemampuan model besar, mendukung ringkasan dokumen, pengenalan layar cerdas, editing gambar AI, dll.}: terhubung dengan HarmonyOS Xiaoyi, mendukung pembuatan copywriting on-device, AI photo editing, AI erase (hapus orang lewat), ringkasan panggilan. Setelah HarmonyOS NEXT 5.0 integrasi tingkat sistem lebih lanjut, kemampuan AI dan sistem melekat dalam.
  \item \textbf{Xiaomi MiLM (Surge Large Model)}: tertanam di HyperOS, mendukung AI portrait, AI ringkasan, AI painter, super resolution. Xiaomi mengambil jalur "kolaborasi device-cloud", model ringan di sisi perangkat, tugas berat di cloud.
  \item \textbf{OPPO (Andes Large Model) dan vivo (Blue Heart Large Model)}: ColorOS 15 dan OriginOS 5 terintegrasi mendalam, mengandalkan AI erase, AI ringkasan dokumen, AI Q\& A gambar, pencarian foto bahasa alami ("cari foto di pantai tahun lalu").
\end{itemize}

\textbf{Saran Pembelian}: jika sering butuh AI untuk memproses dokumen kerja, mengorganisir informasi, utamakan perangkat dengan daya komputasi on-device kuat dan fitur AI terintegrasi tingkat sistem; jika hanya sesekali pakai AI photo editing, ponsel kelas menengah mainstream sudah cukup.

\booksection{Komunikasi Satelit: Dari “Gimmick” Flagship ke Standar Kelas Menengah}

Komunikasi satelit adalah salah satu demokratisasi teknologi paling nyata dua tahun terakhir. Saat debut 2022 Huawei Mate 60 dan iPhone 14 masih eksklusif flagship, 2025 sudah turun ke kelas 2000-3000 yuan.

\begin{itemize}
  \item \textbf{BeiDou short message turun ke ponsel kelas menengah}: fitur short message satelit BeiDou (bisa kirim SMS dalam 1000 karakter Mandarin + lokasi) sudah turun ke model kelas menengah seperti Redmi, realme, iQOO, tanpa biaya bulanan, gratis. Cocok untuk hiking, road trip, dll.
  \item \textbf{Tiantong satellite voice call}: seri Huawei Mate 60/70, seri Honor Magic6 dll. mendukung panggilan suara satelit Tiantong-1, bisa menelpon/menerima panggilan satelit, mengirim pesan satelit. Perlu langganan paket satelit operator (sekitar 10-200 yuan/bulan), cocok untuk pengguna yang sering masuk area terpencil.
  \item \textbf{iPhone satellite messaging (SOS)}: iPhone 14 dan model berikutnya mendukung SOS satelit, bisa mengirim pesan teks ke layanan darurat dan berbagi lokasi. Gratis 2 tahun, sebagian wilayah bisa diperpanjang berbayar.
  \item \textbf{Shooting harian: 1080P/4K + 30fps + encoding H.265, menyeimbangkan kualitas gambar dan ukuran file.}: pengguna kota pada dasarnya tidak butuh, tapi pecinta outdoor, pebisnis yang sering ke daerah terpencil sangat disarankan memilih model dengan fitur satelit — bisa menyelamatkan nyata di momen kritis.
\end{itemize}

\booksection{Tingkat Kematangan Ekosistem Layar Lipat: Bisakah Jadi Ponsel Utama?}

Layar lipat telah bergerak dari “coba teknologi” ke “bisa jadi ponsel utama”. Menilai apakah layar lipat sudah matang, lihat tiga hal: durabilitas hardware, tingkat adaptasi App, ekosistem interaksi.

\begin{itemize}
  \item \textbf{Durabilitas hardware}: umur engsel umumnya mencapai 500 ribu kali+ (buka-tutup 100 kali sehari bisa 13 tahun), engsel waterdrop mengurangi bekas lipatan layar, waterproof IPX8 sudah jadi standar. Layar dalam UTG glass + pelindung membuat layar tidak lagi "rapuh".
  \item \textbf{Adaptasi App}: App utama (WeChat, Taobao, Douyin, Bilibili, Office, WPS) sudah beradaptasi untuk layar besar, tapi masih ada sekitar 30\% App long tail yang belum, tampilannya meregang atau bergaris hitam. Disarankan sebelum membeli coba dulu beberapa App yang sering dipakai.
  \item \textbf{Interaksi mode berdiri}: saat dilipat ke 90°-120°, sebagian App mendukung mode berdiri "layar atas menampilkan konten, layar bawah untuk operasi", cocok untuk video call, selfie, menonton. Kamera selfie mode berdiri adalah pengalaman unik layar lipat.
  \item \textbf{Split-screen dan multi-tasking}: di layar besar bisa membuka dua App split-screen secara bersamaan, bahkan tiga App floating window berdampingan. Logika split-screen HyperOS, ColorOS, OriginOS sudah cukup matang, tapi split-screen iOS (gaya iPadOS) di iPhone belum dibuka.
  \item \textbf{Harga}: lipat horizontal (gaya buku) turun dari level 10 ribu yuan ke kelas 5000-7000 yuan, lipat vertikal (gaya flip) sudah turun ke kelas 4000-5000 yuan. 2025 sebagian pabrikan meluncurkan lipat vertikal kelas 3000 yuan, layar lipat benar-benar masuk pasar massal.
\end{itemize}

\textbf{Saran Pembelian}: jika Anda pembaca berat, pemroses dokumen, pengguna video, layar lipat horizontal layak dipertimbangkan; jika mengejar portabilitas dan atribut sosial, lipat vertikal lebih cocok. Disarankan setidaknya coba di toko offline 30 menit sebelum memutuskan.

\booksection{6G dan Wi-Fi 7: Kecepatan dan Batas Koneksi Generasi Berikutnya}

Meskipun 6G belum benar-benar mendarat, teknologi terkait sudah dalam perjalanan:

\begin{itemize}
  \item \textbf{Perkembangan 6G}: 3GPP merencanakan penyelesaian standar 6G pada 2027, komersial sekitar 2030. Kecepatan puncak teoretis 6G bisa mencapai 1 Tbps (50x 5G), latensi serendah 0.1ms, mendukung "integrasi komunikasi dan penginderaan", ponsel bahkan bisa jadi radar presisi rendah. Saat ini semua pabrikan masih dalam tahap verifikasi prototipe, masih jauh dari produk konsumen.
  \item \textbf{Wi-Fi 7(802.11be)}: sejak 2024 menjadi standar flagship, kecepatan puncak teoretis 46 Gbps (4x Wi-Fi 6), mendukung bandwidth 320MHz, 4K-QAM, MLO (multi-link concurrent). Pengukuran aktual di router Wi-Fi 7, ponsel flagship bisa mencapai 2-5 Gbps, 2-3x lebih cepat dari Wi-Fi 6E. Jika ganti ke ponsel Wi-Fi 7, ingat untuk memasangkan router Wi-Fi 7 agar performa penuh tercapai.
  \item \textbf{Bluetooth 5.4 dan LE Audio}: ponsel baru 2024-2026 umumnya mendukung Bluetooth 5.4 dan LE Audio (audio low power), bisa menyambung dua earphone secara bersamaan untuk dengarkan lagu yang sama, Auracast broadcast audio membuat banyak orang di tempat umum mendengarkan sumber suara yang sama secara sinkron.
\end{itemize}

\booksection{Baterai Silikon-Karbon: Daya Tahan Akhirnya Mendapat Terobosan Level Material}

Daya tahan ponsel selama bertahun-tahun dikritik “cas sehari sekali”, sejak 2024 popularisasi teknologi baterai silikon-karbon memberikan perbaikan substansial.

\begin{itemize}
  \item \textbf{Prinsip teknologi}: kapasitas teoretis anoda grafit tradisional sekitar 372 mAh/g, kapasitas teoretis silikon mencapai 4200 mAh/g (sekitar 10x), tapi silikon mengembang hingga 300\% saat charging-discharging, mudah hancur. Anoda silikon-karbon melalui struktur komposit nano silikon-karbon, mengontrol masalah ekspansi silikon ke rentang yang dapat diterima, meningkatkan densitas energi baterai 15\%-30\%.
  \item \textbf{: Melihat bunga yang tidak dikenal di jalan, buka kamera, AI langsung memberi tahu "ini bunga hibiscus". Mencari foto liburan tahun lalu di pantai di album, cari "laut", AI otomatis menemukan semua foto yang mengandung laut, tanpa perlu mencari manual.}: flagship 2024-2025 yang menggunakan anoda silikon-karbon (seperti Xiaomi 14 Ultra, Honor Magic 6 Ultimate, vivo X100 Pro dll.), dalam bodi 7-8mm memasukkan baterai 5000-6000 mAh, sekitar 800-1000 mAh lebih banyak dari generasi sebelumnya, pemakaian berat bisa bertahan satu setengah hari.
  \item \textbf{Dipadukan dengan pengisian cepat 100W+}: anoda silikon-karbon kompatibel dengan pengisian cepat, pengisian cepat 100W/120W penuh dalam 17 menit bukan lagi berita. 2025 sebagian model mendukung 150W-240W, 10 menit mencapai 80\%.
  \item \textbf{Saran Pembelian}: jika pengguna berat, utamakan model setelah 2024 dengan kombinasi "anoda silikon-karbon + pengisian cepat 100W+", kecemasan baterai akan berkurang drastis.
\end{itemize}

\booksection{Ikhtisar Bab: Satu Kalimat untuk Mengingat Teknologi Terdepan}

\begin{itemize}
  \item Tiga standar ponsel AI: large model on-device, TOPS NPU, AI Agent lintas skenario App.
  \item Komunikasi satelit: BeiDou short message turun ke kelas menengah, Tiantong voice call masih flagship.
  \item Folding screen: hardware matang, App adaptasi 70\%, hover/split screen adalah pengalaman unik.
  \item 6G masih di jalan, Wi-Fi 7 sudah mendarat, Bluetooth LE Audio mulai populer.
  \item Anoda silikon-karbon + pengisian cepat 100W+: terobosan ganda daya tahan dan pengisian.
\end{itemize}

\begin{tcbpractice}{Coba: Buat "daftar kebutuhan teknis" untuk ponsel berikutnya}
Ambil selembar kertas, bagi dalam tiga tingkat “wajib/plus/tidak perlu”, daftar kebutuhan AI, satelit, layar lipat, Wi-Fi 7, pengisian cepat, kapasitas baterai. Saat ke toko offline, “selling point” yang direkomendasikan sales bisa cepat Anda nilai apakah kebutuhan nyata atau sekadar dijual. Misalnya:

\begin{itemize}
  \item Wajib: baterai 5000mAh+, pengisian cepat 80W+, NFC, waterproof IP68.
  \item Plus: AI ringkasan on-device, Wi-Fi 7, BeiDou short message.
  \item Tidak perlu: layar lipat (saya terbiasa candy bar), voice call satelit (tidak bepergian jauh), pengisian cepat 240W ekstrem.
\end{itemize}

\end{tcbpractice}

\begin{tcbsummary}{Ringkasan Bab}
\begin{itemize}
  \item Tiga kriteria ponsel AI: large model on-device, TOPS NPU, skenario AI Agent
  \item Strategi AI pabrikan: Apple fokus privasi dan integrasi, Google on-device open source, Huawei/Xiaomi/OPPO/vivo masing-masing unggul
  \item Komunikasi satelit: BeiDou short message turun ke kelas menengah, Tiantong voice call masih eksklusif flagship
  \item Layar lipat: hardware sudah matang, adaptasi App sekitar 70\%, mode berdiri dan split-screen adalah pengalaman unik
  \item Standar 6G ditetapkan 2027, Wi-Fi 7 sudah populer, Bluetooth LE Audio mulai mendarat
  \item Baterai anoda silikon-karbon membuat 5000-6000mAh menjadi standar flagship, kecemasan baterai berkurang drastis
\end{itemize}

\end{tcbsummary}
